\section{Wstęp}

    \paragraph{}
        Bezzałogowe statki powietrzne stają się jednym z kluczowych elementów współczesnych systemów obserwacyjnych i bojowych.
        Ich zdolność do wielogodzinnego lotu nad trudno dostępnym terenem sprawia,
        że znajdują zastosowanie zarówno w misjach cywilnych,
        jak i operacjach militarnych.

    \paragraph{}
        Rosnące znaczenie tych platform potwierdzają dwa współczesne konflikty zbrojne.
        W wojnie rosyjsko-ukraińskiej drony szczególności FPV z narzędzia rozpoznawczego przekształciły się w niedozowny element uzbrojenia.
        Zyskując na statusie jak legendarny karabinek AK konstrukcji Michaiła Kałasznikowa
        Tworząc rozległe strefy kontroli przestrzeni powietrznej wzdłuż całej linii frontu.
        Ukraina zbudowała wokół nich odrębny rodzaj oddziałów wojskuch np USF Unmanned Systems Forces.
        Z kolei konflikt w Tigraju Etiopia, 2020 2022 pokazuje,
        że dostęp do platform BSP może dosłownie odwrócić losy kampanii: 
        rząd etiopski, 
        stojący w obliczu militarnej klęski.
        Pozyskał chińskie drony Wing Loong 2 oraz tureckie Bayraktar TB2, 
        które zatrzymały natarcie sił Tigraju i umożliwiły kontrofensywę.

    \paragraph{}
        Oba przypadki dowodzą tej samej zasady strona dysponująca stałym rozpoznaniem 
        z powietrza i możliwością precyzyjnego rażenia uzyskuje przewagę operacyjną.
        Której nie równoważy liczebna wyższość wojsk naziemnych, 
        które zreszta wukorzystuja drony tylko tew małe.
        W niniejszej pracy proces obliczeń został zrealizowany przy użyciu języka Python, który pozwolił do bardxziej precyzujnych obliczeń.
        Celem pracy jest opracowanie bezzałogowego systemu klasy MALE o wysokiej autonomii lotu,
        modułowej konstrukcji i możliwości integracji z różnorodnymi sensorami, 
        przeznaczonego do misji rozpoznawczych, 
        patrolowania granic oraz wsparcia działań taktycznych.

    \paragraph{}
        Conflicts such as the ISIS insurgency, the Indo-Pak conflict over Kashmir, Russian
        military actions in Ukraine, the Syrian civil war, the Lebanon conflict, and the tensions between
        the US and North Korea have heightened security concerns. These factors contribute positively
        to the market, thus driving its growth during the forecast period